\documentclass{template-2}

\usepackage{mathtools}
\usepackage{mdframed}
\usepackage{varwidth} 
\usepackage{bbm}
\pgfplotsset{compat=1.18}

\newcommand{\qedwhite}{\hfill \ensuremath{\Box}}

% Random sampling
\newcommand\leftarrowS{\leftarrow\joinrel\smalldollar}
\newcommand\rightarrowS{\smalldollar\joinrel\rightarrow}

\makeatletter
\newcommand{\smalldollar}{\mathrel{\mathpalette\small@dollar\relax}}
\newcommand{\small@dollar}[2]{%
  \vcenter{\hbox{%
    $#1\textnormal{\fontsize{0.7\dimexpr\f@size pt}{0}\selectfont\$}$%
  }}%
}
\makeatother


\begin{document}

\begin{center}
    \Large\textcolor{blue!70!black}{\textbf{Additive Geometry I}} \\
    \large \textbf{Mathematics Seminar at Kyiv School of Economics} \\
    \normalsize\textit{Dmytro Zakharov}
\end{center}


\section{Lattices}\label{section:lattices}

In these notes, we follow~\cite{TaoVu06}.

\begin{definition}
    We define the following notions:
    \begin{enumerate}
        \item[(i)] A \textbf{lattice} $\Lambda \subseteq \mathbb{R}^d$ is any discrete
    additive subgroup of the Euclidean space $\mathbb{R}^d$;
        \item[(ii)] the \textbf{rank} $\text{rk}(\Lambda)$ of $\Lambda$ is the dimension of
    the linear space spanned by elements of $\Lambda$;
        \item[(iii)] if $\text{rk}(L)=d$, we say $L$ has a \textbf{full rank};
        \item[(iv)] if $\Lambda'$ is a lattice and $\Lambda' \subseteq \Lambda \subseteq \mathbb{R}^d$,
        we say $\Lambda'$ is the \textbf{sublattice} of $\Lambda$.
    \end{enumerate}
\end{definition}

\begin{example}
    Here are a couple of examples of lattices:
    \begin{itemize}
        \item[(i)] $\mathbb{Z}^d$ is the full rank lattice of dimension $d$;
        \item[(ii)] $\Lambda := \{\lambda_1(1,2) + \lambda_2(2,1): \lambda_1,
        \lambda_2 \in \mathbb{Z}\} \subseteq \mathbb{R}^2$ is a full rank
        lattice;
        \item[(iii)] In general, for $V = (v_1,\dots,v_k) \in
        \mathbb{R}^{d \times k}$ (matrix of full column rank), the lattice $\Lambda :=
        \mathcal{L}(V) \triangleq \{\sum_{i=1}^k \lambda_iv_i: \forall_{i
        \in [k]} \; \lambda_i \in \mathbb{Z}\}$ is of rank $k$ in $\mathbb{R}^d$;
        \item[(iv)] $q\mathbb{Z}^d$ is a full rank sublattice of $\mathbb{Z}^d$;
        a lattice $\Lambda$ such that $q\mathbb{Z}^d \subseteq \Lambda \subseteq
        \mathbb{Z}^d$ is called a \textbf{$q$-ary lattice}; these are extensively 
        studied in lattice-based cryptography.
    \end{itemize}
\end{example}

We also make a couple of remarks without proof:
\begin{enumerate}
    \item[(a)] if $T: \mathbb{R}^d \to \mathbb{R}^d$ is an invertible linear
    transformation and $\Lambda \subseteq \mathbb{R}^d$ is a lattice, then
    $\Lambda' := T(\Lambda)$ is also a lattice and moreover
    $\text{rk}(\Lambda) = \text{rk}(\Lambda')$;
    \item[(b)] if $\Lambda$ is a lattice, then $\mathbb{R}^d/\Lambda$ is a
    smooth manifold with a natural Lebesgue measure (call it $\mu$) induced
    from $\mathbb{R}^d$;
    \item[(c)] if $\Lambda$ is also of full rank, $\mathbb{R}^d/\Lambda$ is
    compact and thus has a volume $\mu(\mathbb{R}^d/\Lambda)$ which we call
    the \textbf{covolume} of $\Lambda$ (and denote $\text{covol}(\Lambda)$).
\end{enumerate}

Now we wish to classify the lattices. Turns out that any lattice $\Lambda$ can
always be generated by the $\mathbb{Z}$-span of finite collection of linearly
independent vectors. Moreover, one of the vectors $w \in \Lambda$ can be chosen
to be \textit{irreducible}, which means $w/n \not\in\Lambda$ for any $n \in
\mathbb{Z}_{\geq 2}$.

\begin{lemma}[Fundamental theorem of lattices]\label{fundamental-lemma}
    If $\Lambda \subseteq \mathbb{R}^d$ is of rank $k$, then there exist
    linearly independent vectors $v_1,\dots,v_k \in \mathbb{R}^d$ such
    that $\Lambda=\mathcal{L}(v_1,\dots,v_k)$.
\end{lemma}

During construction we will see that $v_1$ can be made
irreducible. One consequence of this lemma is that any lattice of rank $k$ is
isomorphic to $\mathbb{Z}^k$.

\textbf{Proof.} Without loss of generality, assume that $\text{rk}(\Lambda)=d
\geq 1$, else we could pass to a smaller vector space and continue the argument. 

Additionally, observe that $\Lambda$ contains at least one irreducible vector
$w$: indeed, take any non-zero vector $v \in \Lambda$ and take $w$ to be the
smallest of the form $v/n$ (such a vector must exist since $\Lambda$ is
discrete). By the full rank assumption, we find $d$ linearly independent vectors
$v_1,\dots,v_d \in \Lambda$ (with $v_1 := w$), so in particular the volume of
parallelepiped $|\bigwedge_{i=1}^d v_i| > 0$ (\textit{notice that here
$\text{span}(v_1,\dots,v_d)=\mathbb{R}^d$, but $\mathcal{L}(v_1,\dots,v_d) \neq
\Lambda$}). Since $\Lambda_0 := \mathcal{L}(v_1,\dots,v_d) \subseteq \Lambda$ is a sublattice
with covolume $|\bigwedge_{i=1}^d v_i|$, we have $|\bigwedge_{i=1}^d v_i| \geq
\mu(\mathbb{R}^d/\Lambda)$. 

We proceed using the method of descent. We will build the sequence of
sublattices $\{\Lambda_n\}_{n \in \mathbb{Z}}$ (with $\Lambda_n \subseteq
\Lambda$) with the strictly decreasing covolumes (by a factor of $1/2$). Since
we have a positive lower bound, we will conclude that $\Lambda_n = \Lambda$ at
some point.

Here how it goes: consider the fundamental parallelipiped
$\mathcal{P}(\Lambda_0) = \{\sum_{i=1}^d \lambda_iv_i: 0\leq \lambda_i<1\}$.
Notice that it contains some non-zero point $x \in \Lambda$. Thus, write $x =
\sum_{i=1}^d \lambda_iv_i$. Notice that at least one of
$\lambda_2,\dots,\lambda_d$ is non-zero: otherwise, we would have $x=\lambda_1w$
for $\lambda_1 \in [0,1)$, contradicting the fact that $w$ is irreducible. Wlog,
assume that $\lambda_d \neq 0$. 

Now, we claim that we can assume $\lambda_d \leq 1/2$, so that our descent is
fast enough. Indeed, assuming $\lambda_d > 1/2$, we can pick a different point
from $x$: say, $x' = \sum_{i=1}^d v_i - x$, where the corresponding coefficient
is less than $1/2$. Now set $\Lambda_1 := \mathcal{L}(v_1,\dots,v_{d-1},x)$.
Notice that $\text{covol}(\Lambda_1) \leq \frac{1}{2}\text{covol}(\Lambda_0)$.

Define $\Lambda_n$ from $\Lambda_{n-1}$ recursively using the same steps. At
some point we terminate and thus get the required $\mathbb{Z}$-basis.

\begin{lemma}[Corollary]\label{fundamental-corollary}
    This lemma has a couple of almost direct corollaries:
    \begin{enumerate}
        \item[(a)] \textbf{Splitting lemma.} Let $\Lambda$ be a lattice of rank $k$,
        and let $v \in \Lambda$ be an irreducible vector. Then, there exists a
        sublattice $\Lambda' \subseteq \Lambda$ of rank $k-1$ such that $\Lambda
        = \Lambda' + \mathbb{Z}v$ with $\Lambda' \cap \mathbb{Z}v = \emptyset$.
        \item[(b)] \textbf{Fundamental theorem of f.g. torsion-free additive groups.}
        Let $Z$ be a finitely generated torsion-free additive group. Then, $Z
        \cong \mathbb{Z}^d$ for some $d \geq 0$.
    \end{enumerate}
\end{lemma}

\textbf{Proof.}

\textit{Part (a).} Apply Lemma~\ref{fundamental-lemma} with $v_1 := v$, and 
set $\Lambda' := \mathcal{L}(v_2,\dots,v_k)$.

\textit{Part (b).} Since $Z$ is finitely generated, let $Z=(z_1,\dots,z_n)$ (for
$z_i \in Z$). Define the lattice $\Lambda := \{u \in \mathbb{Z}^n: \sum_{i=1}^n
u_iz_i = 0\}$. Now notice that it is precisely the kernel of the homomorphism
$(u_1,\dots,u_n) \mapsto \sum_{i=1}^n u_iz_i$. Therefore, it induces the
isomorphism $Z \cong \mathbb{Z}^n/\Lambda$. In particular,
$\mathbb{Z}^n/\Lambda$ is torsion-free. 

Now we proceed by induction. Base case is trivial and assume $\Lambda \neq
\{0\}$. Pick some irreducible vector $v \in \Lambda$. Note that it is also
irreducible in $\mathbb{Z}^n$: if $v=mw$ for some $w \in \mathbb{Z}^n$ and
$m>1$, then $w+\Lambda$ would be a non-zero element of $\mathbb{Z}^n/\Lambda$
such that $m(w+\Lambda)=0+\Lambda$, contradicting the torsion-free assumption
for $\mathbb{Z}^n/\Lambda$. This implies that $\mathbb{Z}^n/v\mathbb{Z} \cong
\mathbb{Z}^{n-1}$. We finish the proof by noticing that
\begin{equation*}
    \mathbb{Z}^n/\Lambda \cong \frac{\mathbb{Z}^n/(v\mathbb{Z})}{\Lambda/(v\mathbb{Z})} \cong \mathbb{Z}^{n-1}/\Lambda'.
\end{equation*}

\section{Quotients of lattices}

As already demonstrated in Lemma~\ref{fundamental-corollary} (\textit{part
(b)}), quotiening two lattices is a powerful tool. More generally, if $G$ is a
finitely generated additive group by $d$ elements $v_1,\dots,v_d \in G$, then $G
\cong \mathbb{Z}^d/\Lambda_G$ where $\Lambda_G := \{u \in \mathbb{Z}^d: \langle
u, v \rangle = 0\}$. Therefore, we provide the following tool to help analyzing
quotients of two lattices.

\begin{theorem}[Smith normal form]
    Let $\Lambda,\Lambda' \subseteq \mathbb{R}^d$ be two full-rank lattices and
    $\Lambda' \subseteq \Lambda$. Then, there exist linearly independent
    $v_1,\dots,v_d \in \Lambda$ such that
    \begin{equation*}
        \Lambda = \mathcal{L}(v_1,\dots,v_d), \quad \Lambda' = \mathcal{L}(N_1v_1,\dots,N_dv_d),
    \end{equation*}
    with $N_1 \geq 1$ and $N_j \mid N_{j+1}$ for all $1 \leq j < d$.
\end{theorem}

\textbf{Proof.} We induce on $d$ with $d=0$ being vacuously true. For any
non-zero $v \in \Lambda$ define the \textit{index} of $v$ (call it
$\text{ind}_{\Lambda}(v)$) as follows:
\begin{equation*}
    \text{ind}_{\Lambda}(v) = \max \{n \in \mathbb{N}: v/n \in \Lambda\}.
\end{equation*}

Note that since $\Lambda$ is discrete, $\text{ind}_{\Lambda}(v) < \infty$. Also
note that $\text{ind}_{\Lambda}(v) = n$ iff $v=nw$ for irreducible $w \in \Lambda$.

Since $\Lambda'$ has full rank, it contains non-zero vectors, each of which has
an index. Let $N_1 := \min_{v \in \Lambda'} \text{ind}_{\Lambda}(v)$. By the
well-ordering principle there exists an irreducible vector $v_1 \in \Lambda$
such that $N_1v_1 \in \Lambda'$.

Using Lemma~\ref{fundamental-lemma}, we may apply an invertible transformation
to map $\Lambda$ to $\mathbb{Z}^d$ such that $v_1$ gets mapped to the standard
basis vector $e_1$. Now let $(n_1,\dots,n_d)$ be any vector in $\Lambda'$.
Observe that $n_1,\dots,n_d$ are integers; furthermore, $n_1$ must be a multiple
of $N_1$. Thus, we may factorize $\Lambda' = (N_1e_1)\mathbb{Z} + \Lambda''$
where $\Lambda'' \subseteq \mathbb{Z}^{d-1}$ is a lattice. Also note that if $x
\in \Lambda''$, then $(N_1,x) \in \Lambda'$ and hence (since $(N_1,x)$ must have
an index of at least $N_1$), $x$ must be a multiple of $N_1$. Thus, $\Lambda''
\subseteq N_1 \mathbb{Z}^{d-1}$, so we may write $\Lambda' =
N_1(e_1\mathbb{Z}+\Lambda''')$ for some sublattice $\Lambda''' \subseteq
\mathbb{Z}^{d-1}$. Clearly $\text{rk}(\Lambda''')=d-1$.

We now invoke the inductive hypothesis. By applying an invertible transformation
to $\mathbb{Z}^{d-1}$ if necessary, we may assume that
\begin{equation*}
    \Lambda''' = \{(n_2M_2,\dots,n_dM_d): n_2,\dots,n_d \in \mathbb{Z}\}
\end{equation*}

for some $1 \leq M_2 \leq \dots \leq M_d$ such that $M_j \mid M_{j+1}$ for every
$j$. The claim follows by setting $N_j := N_1M_j$ for $j>1$.

\begin{lemma}[Corollary]\label{smith-corollary}
    This lemma also has two nice corollaries:
    \begin{enumerate}
        \item[(a)] \textbf{Fundamental theorem of finite additive groups.} Every 
        finite additive group $G$ is isomorphic to the direct sum of a finite 
        number of cyclic groups $\mathbb{Z}/N\mathbb{Z}$.
        \item[(b)] \textbf{Fundamental theorem of f.g. additive groups.}
        Every finitely generated additive group $G$ is isomorphic to the direct 
        sum of a finite number of cyclic groups $\mathbb{Z}/N\mathbb{Z}$, and 
        a lattice $\mathbb{Z}^d$ for some $d \geq 0$.
    \end{enumerate}
\end{lemma}

\textbf{Proof.} 

\textit{Part (a).} Since the group is finitely generated, let
$G=(g_1,\dots,g_d)$. Define the map $\phi: \mathbb{Z}^d \to G$ as
$\phi(n_1,\dots,n_d) = \sum_{i=1}^d n_ig_i$. Using first isomorphism theorem, we
as before note that $G \cong \mathbb{Z}^d/\ker(\phi) \leq
\mathbb{R}^d/\ker(\phi)$. Clearly, $\ker(\phi)$ is a lattice of some rank $0
\leq k \leq d$. Moreover, it must be of full rank, otherwise
$\mathbb{Z}^d/\ker(\phi)$ would be infinite, and thus so would hold for $G$.
Hence by Lemma~\ref{fundamental-lemma} it is generated by $d$ linearly independent
vectors $v_1,\dots,v_d \in \mathbb{Z}^d$. Using the Smith normal 
form, we can after applying an isomorphism write $\ker(\phi)$ as the 
lattice generated by $N_1e_1,\dots,N_de_d$ for some $N_1,\dots,N_d \geq 1$;
this makes $G$ isomorphic to $G \cong \bigoplus_{i=1}^d (\mathbb{Z}/N_i\mathbb{Z})$
as desired (moreover, $N_i \mid N_{i+1}$ for each $i \geq 1$).

\textit{Part (b).} Let $\text{tor}(G) := \{x \in G: \exists_{n>0} \; nx=0\}$ be
the torsion group of $G$. By \textit{part (a)} we know that $\text{tor}(G)$ is
the direct sum of cyclic groups. The quotient group $G/\text{tor}(G)$ is
torsion-free and is thus isomorphic to $\mathbb{Z}^d$ for some $d \geq 0$ by
Corollary~\ref{fundamental-corollary}. Let $\widehat{e}_1,\dots,\widehat{e}_d$
be arbitrary representatives in $G$ of the standard basis $e_1,\dots,e_d$ of
$\mathbb{Z}^d$. Then, $G \cong \text{tor}(G) \oplus \bigoplus_{i=1}^d
\mathbb{Z}\widehat{e}_i \cong \text{tor}(G) \oplus \mathbb{Z}^d$.

\section{Convex bodies}

Denote by $\mu(A)$ the volume of a set $A \subseteq \mathbb{R}^d$. If $A \in
\mathbb{R}^d$ and $\lambda \in \mathbb{R}$, denote $\lambda \cdot A \triangleq
\{\lambda a: a \in A\}$. Observe that $\mu(\lambda \cdot A)=|\lambda|^d\mu(A)$.

\begin{definition}\label{def:convex-set}
    We give a couple of new definitions:
    \begin{enumerate}
        \item[(a)] a set $A \subseteq \mathbb{R}^d$ is
        \textbf{convex} if $(1-\theta)x + \theta y \in A$ whenever $x,y \in A$ and
        $\theta \in [0,1]$;
        \item[(b)] a set $A$ is called a \textbf{convex body} if it is convex,
        open, non-empty, and bounded;
        \item[(c)] a convex body $A$ is \textbf{symmetric} if $A=-A$.
    \end{enumerate} 
\end{definition}

\begin{remark}
    One can show that Definition~\ref{def:convex-set} is equivalent to the
    following: the set $A$ is convex iff $\lambda \cdot A + \mu \cdot B =
    (\lambda+\mu)\cdot A$ for all $\lambda,\mu \geq 0$. In particular, for
    convex sets we have $nA = |n|\cdot A$ for any $n \in \mathbb{Z}$.
\end{remark}

Notice that $\mu(A+A)=\mu(2 \cdot A)=2^d\mu(A)$. For $A-A$, we can use
\begin{lemma}
    For any bounded open subsets $X,Y,Z \subseteq \mathbb{R}^d$ (not necessarily 
    convex):
    \begin{equation*}
        \mu(X-Z)\mu(Y) \leq \mu(X-Y)\mu(Y-Z).
    \end{equation*}
\end{lemma}

Using this lemma (set $X:=A$, $Y:=-A$, $Z:=A$), we have $\mu(A-A) \leq
4^d\mu(A)$. In the converse direction, Brunn-Minkowski inequality gives
$\mu(A-A) \geq 2^d\mu(A)$ (which we will hopefully discuss in the future).

\begin{theorem}[John's Theorem]
    Let $A \subseteq \mathbb{R}^d$ be a convex body. Then there exists an
    invertible linear map $T: \mathbb{R}^d \to \mathbb{R}^d$ and a point $x_0
    \in A$ such that $B_d \subseteq T(A-x_0) \subseteq d \cdot B_d$, where $B_d
    = \{x \in \mathbb{R}^d: \|x\|_2 < 1\}$. If $A$ is symmetric, then we can
    improve these inclusions to $B_d \subseteq T(A) \subseteq \sqrt{d} \cdot
    B_d$.
\end{theorem}

\textbf{Proof sketch.} Define the \textit{ellipsoid} to be any set $E$ of the form
$E=L(B_d)+x_0$ where $x_0 \in \mathbb{R}^d$ and $L$ any linear map (not
necessarily invertible). 

Since $A$ is open and bounded, it is easy to see that the set of all ellipsoids
$E$ contained in $A$ is a compact set. Besides, the volume of ellipsoid $E$ is
$\mu(E)=|\det L|$, which is clearly a continuous function of $E$. Thus there
exists an ellipsoid $E=L(B_d)+x_0$ in $A$ that maximizes the volume $\mu(E)$;
since $A$ is open, this volume is non-zero, and hence $L$ is invertible. By
applying $L^{-1}$, we may thus assume that $E$ is a translate $E=B_d+y_0$ of a
unit ball, where $y_0=L^{-1}(x_0)$. Wlog, now assume that we have such an $E$ in
$A$, keeping in mind that $T=L^{-1}$.

Assume for now that $A$ is symmetric. Observe that since $A$ contains $B_d+y_0$,
then it also contains $B_d-y_0$ by symmetry, and hence contains $B_d$ since it
is in the convex hull of $B_d+y_0$ and $B_d-y_0$. To conclude the proof of the
lemma in this case we need to show that $A \subseteq \sqrt{d} \cdot B_d$. 

Suppose for contradiction this is not the case, so suppose wlog that $re_1 \in
A$ for some $r>\sqrt{d}$ and $e_1$ the first basis vector. The rest of the proof
consists in constructing some ellipsoid $L$ contained in $A$ but with
$\mu(L)>\mu(B_d)$ (contradicting the minimality of $B_d$). Specifically,
\begin{equation*}
    L := L_{\varepsilon,\delta}(B_d), \quad L_{\varepsilon,\delta} := \begin{pmatrix}
        1+(\sqrt{d}-1+\delta)\varepsilon & 0 \\
        0 & (1-\varepsilon)\mathbb{I}_{d-1}
    \end{pmatrix}.
\end{equation*}

A couple of statements that finishes the proof:
\begin{enumerate}
    \item $L_{0,\delta}(B_d)=B_d$ (that is, for $\varepsilon=0$);
    \item $\det
    L_{\varepsilon,\delta}=(1-\varepsilon)^{d-1}(1+(\sqrt{d}-1+\delta)\varepsilon)$
    has a positive $\varepsilon$-derivative at $\varepsilon=0$;
    \item for small enough $\delta$, we have $L_{\varepsilon,\delta}(B_d) \subseteq A$.
\end{enumerate}

For symmetric case turns out that we can use the same transformation
$L_{\varepsilon,\delta}$, but to ensure that the third condition is met, 
we need the radius of $d$ for the ball.

\subsection*{References}
\footnotesize
\renewcommand{\refname}{\vspace{-1em}}
\bibliographystyle{alphaurl}
\bibliography{refs-lattices}

\end{document}